\subsection{Примеры, основные понятия}

\subsubsection*{Откуда берутся интегральные уравнения}

Все четыре примера ниже — это одна и та же мысль: задача для обыкновенного
дифференциального уравнения вместе с начальными или краевыми условиями
переписывается в виде одного интегрального уравнения, в которое условия уже
встроены.

\begin{example}[задача Коши первого порядка]
\[
  \begin{cases}
    u' = f(x,u), \\
    u(x_0) = u_0 .
  \end{cases}
\]
Проинтегрируем уравнение от $x_0$ до $x$:
\[
  u(x) - u_0 = \int\limits_{x_0}^{x} f\big(y, u(y)\big)\, dy
  \qquad\Longrightarrow\qquad
  u(x) = \int\limits_{x_0}^{x} f\big(y, u(y)\big)\, dy + u_0 .
\]

\medskip
В частном случае линейного уравнения
\[
  u' + p(x)\,u = q(x), \qquad\text{то есть}\qquad f(x,u) = -p(x)\,u + q(x) ,
\]
получаем линейное интегральное уравнение
\[
  u(x) + \int\limits_{x_0}^{x} p(y)\, u(y)\, dy
  = u_0 + \int\limits_{x_0}^{x} q(y)\, dy .
\]
\end{example}

\begin{remark}
В операторной записи это $u = Au + a$ — простейший вид, к которому сводится
дело. Отсюда же напрашивается и способ решения: итерации
\[
  u_n = A u_{n-1} + a .
\]
\end{remark}

\begin{example}[задача Коши второго порядка]
\[
  \begin{cases}
    u'' = f(x,u), \\
    u(x_0) = u_0, \quad u'(x_0) = u_1 .
  \end{cases}
\]
Первое интегрирование даёт
\[
  u'(x) - u_1 = \int\limits_{x_0}^{x} f\big(y,u(y)\big)\, dy ,
\]
второе —
\[
  u(x) = \int\limits_{x_0}^{x} (x-y)\, f\big(y,u(y)\big)\, dy
       + u_1\,(x - x_0) + u_0 .
\]
\end{example}

\begin{idea}
Множитель $(x-y)$ — это формула Коши для повторного интегрирования: двойной
интеграл $\int_{x_0}^{x}\!\int_{x_0}^{s} g(y)\,dy\,ds$ после перемены порядка
интегрирования превращается в однократный $\int_{x_0}^{x}(x-y)g(y)\,dy$. При
$n$-кратном интегрировании множитель был бы
$\frac{(x-y)^{n-1}}{(n-1)!}$.
\end{idea}

\begin{example}[краевая задача]\label{ex:dirichlet-kernel}
\[
  \begin{cases}
    u'' = f(x,u), & x \in (0,\ell), \\
    u(0) = 0, \quad u(\ell) = 0 .
  \end{cases}
\]
Общее решение берём из предыдущего примера, но с произвольными постоянными:
\[
  u(x) = \int\limits_{0}^{x} (x-y)\, f\big(y,u(y)\big)\, dy + c_1 x + c_2 .
\]
Условие $u(0) = 0$ даёт $c_2 = 0$. Условие $u(\ell) = 0$ даёт
\[
  0 = \int\limits_{0}^{\ell} (\ell-y)\, f\big(y,u(y)\big)\, dy + c_1 \ell
  \qquad\Longrightarrow\qquad
  c_1 = -\frac1\ell \int\limits_{0}^{\ell} (\ell-y)\, f\big(y,u(y)\big)\, dy .
\]
Подставляя,
\[
  u(x) = \int\limits_{0}^{x} (x-y)\, f\big(y,u(y)\big)\, dy
       - \frac{x}{\ell}\int\limits_{0}^{\ell} (\ell-y)\, f\big(y,u(y)\big)\, dy .
\]
\end{example}

Приведём это к одному интегралу по $[0,\ell]$. Разобьём второй интеграл на
$\int_0^x$ и $\int_x^\ell$ и соберём коэффициенты при $f(y,u(y))$:
\[
  y \le x: \quad
  (x-y) - \frac{x}{\ell}(\ell - y) = -y + \frac{x}{\ell}\,y
         = -\frac{y}{\ell}\,(\ell - x) ,
\]
\[
  y > x: \quad
  -\frac{x}{\ell}(\ell - y) = -\frac{x(\ell-y)}{\ell} .
\]
Значит
\[
  \boxed{\ u(x) + \int\limits_{0}^{\ell} K(x,y)\, f\big(y,u(y)\big)\, dy = 0\ }
\]
где
\[
  K(x,y) =
  \begin{cases}
    \dfrac{x\,(\ell - y)}{\ell}, & 0 \le x \le y \le \ell, \\[10pt]
    \dfrac{y\,(\ell - x)}{\ell}, & 0 \le y \le x \le \ell,
  \end{cases}
  \qquad K(x,y) = K(y,x), \qquad K \ge 0 .
\]

\inputfig{figures/fig-kernels}

\begin{idea}
$K(x,y)$ — это \textbf{функция Грина} задачи Дирихле для оператора
$\frac{d^2}{dx^2}$ на отрезке $[0,\ell]$. Её симметричность $K(x,y)=K(y,x)$ не
случайна: она отражает самосопряжённость оператора с этими краевыми условиями.
Неотрицательность $K \ge 0$ — это принцип максимума в простейшем виде.

\medskip
Принципиальная разница с примерами~1 и~2: там верхний предел интегрирования
был переменным ($x$), здесь он фиксирован ($\ell$). Это и есть разница между
уравнениями Вольтерра и Фредгольма — см. ниже.
\end{idea}

\subsubsection*{Классификация}

\begin{example}[общий вид]
Пусть $D$ — область, $K(x,y)$ — заданная функция двух переменных,
$f(x)$ — заданная функция. Рассматриваются уравнения относительно $u$:
\[
  \text{I род:}\qquad
  \int\limits_{D} K(x,y)\, u(y)\, dy = f(x), \qquad x \in D ,
\]
\[
  \text{II род:}\qquad
  u(x) - \int\limits_{D} K(x,y)\, u(y)\, dy = f(x), \qquad x \in D .
\]
Функция $K(x,y)$ называется \textbf{ядром}. Это уравнения
\textbf{Фредгольма} первого и второго рода соответственно.
\end{example}

\begin{definition}
Уравнения с переменным верхним пределом
\[
  \int\limits_{a}^{x} K(x,y)\, u(y)\, dy = f(x), \qquad x \ge a ,
\]
\[
  u(x) - \int\limits_{a}^{x} K(x,y)\, u(y)\, dy = f(x), \qquad x \ge a
\]
называются уравнениями \textbf{Вольтерра} первого и второго рода.
\end{definition}

\begin{remark}
Уравнение Вольтерра — частный случай уравнения Фредгольма: достаточно
положить
\[
  K(x,y) = 0 \quad\text{при } y > x ,
\]
и интеграл по всей области автоматически обрежется до $\int_a^x$
(рис.~\ref{fig:kernels}\,б).
\end{remark}

\begin{idea}
Разница между родами куда существеннее, чем между Фредгольмом и Вольтерра. В
уравнении \textbf{второго} рода искомая функция входит и вне интеграла — такое
уравнение по структуре близко к $u - Au = f$, то есть к обращению оператора
$I - A$, и решается итерациями. В уравнении \textbf{первого} рода $u$ есть
только под интегралом; интегрирование «сглаживает», поэтому обратная операция
— дифференцирование — неустойчива. Уравнения первого рода, как правило,
некорректны по Адамару.
\end{idea}

\subsubsection*{Интегральный оператор}

Чаще уравнение второго рода пишут с параметром:

\stepcounter{equation}
\begin{equation}\tag{1}\label{eq:fredholm2}
  u(x) - \lambda \int\limits_{D} K(x,y)\, u(y)\, dy = f(x),
  \qquad x \in D ,
\end{equation}
при $\lambda = 1$ получается форма, выписанная выше.

\begin{definition}
Отображение
\[
  (Au)(x) = \int\limits_{D} K(x,y)\, u(y)\, dy
\]
называется \textbf{интегральным оператором} с ядром $K$. Он линеен, и
уравнение~\eqref{eq:fredholm2} записывается как
\[
  u - \lambda A u = f .
\]
\end{definition}

\begin{idea}
Параметр $\lambda$ введён не для красоты: уравнение $u - \lambda Au = f$ —
это задача об обращении оператора $I - \lambda A$, и ответ на неё зависит от
того, является ли $1/\lambda$ собственным значением $A$. Отсюда и вся теория
Фредгольма, к которой курс дальше и идёт.
\end{idea}
